\documentclass[11pt,a4paper]{article}
\usepackage[margin=25mm]{geometry}
\usepackage[T1]{fontenc}
\usepackage{lmodern,amsmath,amssymb,amsthm,mathtools,booktabs,microtype}
\usepackage[hidelinks]{hyperref}
\newtheorem{theorem}{Theorem}
\newtheorem{lemma}[theorem]{Lemma}
\newtheorem{proposition}[theorem]{Proposition}
\newtheorem{corollary}[theorem]{Corollary}
\theoremstyle{remark}
\newtheorem{remark}[theorem]{Remark}
\newcommand{\ind}{\mathbf 1}
\newcommand{\EE}{\mathbb E}
\newcommand{\PP}{\mathbb P}
\newcommand{\FF}{\mathcal F}
\title{Arithmetic Rigidity of Late $3x+c$ Returns\\
\large Deterministic parity-strip constraints and the entropy barrier at a stopped entrance}
\author{Independent research note --- not inserted into the main manuscript}
\date{16 November 2026}
\begin{document}
\maketitle
\begin{abstract}
Starting from the dyadic block-leakage reduction for
actual stopped Syracuse trajectories, we isolate an elementary but useful
\emph{unconditional dynamical constraint}: every return to the same
multiplicative shell before a lower first-passage cutoff forces the
cumulative accelerated-map parity count to be within a narrow strip about
slope $\log 2/\log 3\approx0.63093$.  This holds for arbitrary late
return times, not just the classical finite coupling horizon.  We prove
an explicit finite-horizon counting inequality and, equally importantly,
a sharp obstruction to extending its Bernoulli estimate past $O(\log X)$
using only initial residue-class equidistribution.  These results do
\emph{not} prove polynomial survival of the stopped first-excess law.
We give a precise two-part target for subsequent genuinely new work.
\end{abstract}

\section{The unrestricted-time arithmetic identity}
Let $c$ be an odd integer with $c\ge-1$, and write
\[
 T_c(n)=\begin{cases}n/2,&n\equiv0\pmod2,\\(3n+c)/2,&n\equiv1\pmod2.\end{cases}
\]
For $n_i=T_c^i(n)$, define the count of odd steps
$H_t(n)=\sum_{i=0}^{t-1}\ind_{\{n_i\equiv1\pmod2\}}$ and
$\rho=\log 2/\log 3$.
Fix a cutoff $a\ge\max(1,2|c|)$; all the formulas below involve
integer states strictly larger than $a$.

\begin{lemma}[Exact log-cocycle and uniform cutoff error]\label{lem:cocycle}
If $n_0,\ldots,n_{t-1}>a$, then
\begin{equation}\label{eq:exact}
 \log\frac{n_t}{n_0}=H_t(n)\log3-t\log2+
 \sum_{\substack{0\le i<t\\ n_i\ \mathrm{odd}}}
 \log\!\left(1+\frac{c}{3n_i}\right).
\end{equation}
In particular, writing
$\epsilon_a=|c|/((3a-|c|)\log3)$,
\begin{equation}\label{eq:drift}
\left|H_t(n)-\rho t-\frac{\log(n_t/n_0)}{\log3}\right|
\le \epsilon_a t.
\end{equation}
\end{lemma}
\begin{proof}
At an even step $n_{i+1}/n_i=1/2$.  At an odd step
$n_{i+1}/n_i=(3/2)(1+c/(3n_i))$. Multiplication and taking
logarithms gives \eqref{eq:exact}. Since $|c|/(3n_i)<|c|/(3a)<1$,
the elementary bound $|\log(1+u)|\le |u|/(1-|u|)$ gives
$|\log(1+c/(3n_i))|\le |c|/(3a-|c|)$.
Summing at most $t$ terms and dividing by $\log3$ proves
\eqref{eq:drift}.
\end{proof}

\begin{theorem}[Every return must revisit a high-parity strip]\label{thm:strip}
Let $I=(Y/2,Y]\cap\mathbb N$ be any multiplicative shell with
$Y/2>a$.  If $n_0,n_t\in I$ and $n_i>a$ for every $0\le i<t$, then
\begin{equation}\label{eq:strip}
\boxed{\quad |H_t(n_0)-\rho t|\le
C_0+\epsilon_a t,\qquad C_0=\frac{\log2}{\log3}.\quad}
\end{equation}
Thus for any $\theta\in(1/2,\rho)$ there are explicit
$a_0(c,\theta)$ and $t_0(\theta)$ such that, whenever $a\ge a_0$
and $t\ge t_0$, every such return satisfies
\begin{equation}\label{eq:biased}
\boxed{\quad H_t(n_0)\ge \theta t.\quad}
\end{equation}
One admissible choice is
$\epsilon_a\le(\rho-\theta)/2$ and
$t_0=\left\lceil2C_0/(\rho-\theta)\right\rceil$.
The assertion is completely deterministic and makes no assumption on
the distribution of parity steps or the size of $t$.
\end{theorem}
\begin{proof}
As both endpoints lie in $I$, their ratio is in $(1/2,2)$, so
$|\log(n_t/n_0)|\le\log2$.  Apply Lemma~\ref{lem:cocycle}.
The final assertion follows from
$H_t/t\ge \rho-C_0/t-\epsilon_a\ge\theta$.
\end{proof}

\begin{remark}[Why this applies beyond the logarithmic horizon]
As $a_X=X^{1-2\beta}\to\infty$ in the main manuscript, the error
$\epsilon_{a_X}t$ in \eqref{eq:strip} remains $o(1)$ for
$t=o(a_X)$, including polynomially long times such as
$t\le a_X^{1/2}$.  This is a \emph{pathwise} relation, not a
probabilistic estimate for how often such paths occur.
All step counts here refer to the accelerated $T_c$ map; the odd-to-odd
Syracuse map groups several of these steps into one block.
\end{remark}

\section{An explicit arithmetic counting bound}
For a probability measure $\nu$ on a finite set of positive integers,
let
\[
 d_\nu(t)=2^t\max_{r\in\mathbb Z/2^t\mathbb Z}
  \nu\{n:n\equiv r\pmod{2^t}\},\qquad t\ge1.
\]
This is the exact maximal density of the first $t$ parity bits
relative to their uniform Bernoulli distribution.  Let
$D(\theta\|1/2)=\theta\log(2\theta)
+(1-\theta)\log(2(1-\theta))$.

\begin{lemma}[Parity words are a residue permutation]\label{lem:parity}
For every odd $c$ and $t\ge1$, the map
\[
 n\bmod2^t\longmapsto
 (n_0\bmod2,\ldots,n_{t-1}\bmod2)
\]
is a bijection between $\mathbb Z/2^t\mathbb Z$ and the binary words
of length $t$.
\end{lemma}
\begin{proof}
Induct on $t$.  Fixing the first parity bit writes $n=2v$ or
$n=2v+1$.  The first iterate is respectively $v$ or
$3v+(3+c)/2$, each a bijection of the residue classes
$v\bmod2^{t-1}$.  The induction hypothesis uniquely determines
those residue classes from the remaining $t-1$ bits.
\end{proof}

\begin{proposition}[Stopped-return counting at an arbitrary time]\label{prop:count}
Choose $\theta\in(1/2,\rho)$, $a$ and $t_0$ as in
Theorem~\ref{thm:strip}, and let $\nu$ be any probability measure
supported on a shell $I\subset(a,\infty)$.  For every $t\ge t_0$,
\begin{equation}\label{eq:count}
\nu\{n:T_c^t(n)\in I,\;T_c^i(n)>a\ (0\le i<t)\}
\le d_\nu(t)\exp[-tD(\theta\|1/2)].
\end{equation}
Writing $N_{I,T}(n)$ for the number of visits to $I$ at accelerated
times $0\le t\le T$ before first passage at or below $a$, one also
has, for $k-1\ge t_0$,
\begin{equation}\label{eq:count-visits}
\nu\{N_{I,T}\ge k\}\le
\sum_{t=k-1}^{T}d_\nu(t)\exp[-tD(\theta\|1/2)].
\end{equation}
\end{proposition}
\begin{proof}
Theorem~\ref{thm:strip} places the return event inside
$\{H_t\ge\theta t\}$.  Lemma~\ref{lem:parity} shows that the latter
comprises $\sum_{h\ge\theta t}\binom{t}{h}$ residue classes
modulo $2^t$.  Each such class has $\nu$-mass at most
$2^{-t}d_\nu(t)$, and the standard binomial Chernoff bound gives
\[
 2^{-t}\sum_{h\ge\theta t}\binom{t}{h}
 \le e^{-tD(\theta\|1/2)}.
\]
If there are at least $k$ visits by time $T$, starting with the
initial one, at least one further visit occurs at some $t\ge k-1$.
The union bound proves \eqref{eq:count-visits}.
\end{proof}

\begin{proposition}[An entropy barrier for finite stopped entrance laws]
\label{prop:barrier}
If $\nu$ is supported on at most $M$ positive integers, then
\begin{equation}\label{eq:entropy-floor}
\boxed{\quad d_\nu(t)\ge\frac{2^t}{M}\quad\text{for every }t.\quad}
\end{equation}
Consequently no $X$-uniform estimate $d_{\nu_X}(t)\le
C\exp(\zeta t)$ with $\zeta<\log2$ can hold for
arbitrarily large $t/\log M_X$, when $\nu_X$ is supported on
$M_X$ finite states.  In particular, for the first-excess entrance
law on $W_X$ one has $M_X\le X+1$, so a Bernoulli Chernoff bound based
\emph{only on the residue distribution at the initial stopped state}
can never yield a decaying estimate throughout the polynomial-time
range $t\sim a_X^{1/2}$.
\end{proposition}
\begin{proof}
The largest atom of a probability measure on at most $M$ points is
at least $1/M$.  Its residue class has at least that mass, whence
\eqref{eq:entropy-floor}.  The claimed incompatibility follows by
comparing $e^{t\log2}/M_X$ with $Ce^{\zeta t}$ as
$t/\log M_X\to\infty$.  Note that
$D(\theta\|1/2)<\log2$, so \eqref{eq:count} indeed loses power
in this range without \emph{new dynamical} information.
\end{proof}

\section{What actually remains for first-excess block leakage}
Let $\eta_X$ be the actual subprobability distribution of the
first-excess state in $W_X$, as defined in the preceding independent
research note.  Its mass $p_X\to0$ is covered by the manuscript's
tightness estimate.  For a shell $I_{X,j}$, let $E_{X,j}$ denote the
exact first-future-entry kernel and let $z$ be the first such state.
Write $\tau_a(z)$ for the first passage to $[1,a]$.

For $T\ge1$, define the arithmetic strip local time
\[
 \mathcal L_{a,T}(z)=
 \#\bigl\{1\le t<\min(T,\tau_a(z)):
 |H_t(z)-\rho t|\le C_0+\epsilon_a t\bigr\}.
\]
When the entry shell lies above $a$, every visit to the shell by
time $T$ lies inside this strip, by Theorem~\ref{thm:strip}.

\begin{proposition}[Exact two-obstruction decomposition]
\label{prop:two-obstructions}
For any integer $n\ge2$ and horizon $T\ge1$, the actual return
survival mass defined in the dyadic block-leakage note obeys
\begin{align}\label{eq:two-obstructions}
 q_{X,j}(n)
 &\le
 \eta_X E_{X,j}\ind_{\{\mathcal L_{a_X,T}\ge n-1\}}
 +\eta_X E_{X,j}\ind_{\{\tau_{a_X}>T\}}.
\end{align}
Thus to prove an actual polynomial return bound by this route,
one needs to control \emph{both} the stopped-law parity-strip local
time and the mass of orbits surviving beyond the mesoscopic horizon.
Neither bound follows from the original finite-cylinder coupling.
\end{proposition}
\begin{proof}
If a first shell entry $z$ has at least $n$ shell visits before the
cutoff and its orbit reaches the cutoff by time $T$, then at least
$n-1$ of these visits occur at positive times strictly below
$\tau_{a_X}(z)\le T$.  Every such time is counted by
$\mathcal L_{a_X,T}(z)$, so its local time is at least $n-1$.
Otherwise $\tau_{a_X}(z)>T$.  Sum these two exhaustive alternatives
against the actual entrance measure $\eta_X E_{X,j}$.
\end{proof}

\begin{remark}[What is and is not new]
The parity-word bijection and the binomial large-deviation estimate
are classical (Terras); they are not offered as new contributions.
The deterministic cocycle is elementary.  The useful distinction in
this note is that \emph{pathwise} return rigidity persists at very
late times, whereas \emph{distributional} randomness provably cannot
be inferred from finite-support residue equidistribution there.
Proposition~\ref{prop:two-obstructions} identifies two distinct tasks
for an actual-orbit proof.  Establishing either task uniformly for
the stopped first-excess law remains open.  We do not claim
$T_X=o(X)$, a polynomial return theorem, or a solution to the
Collatz conjecture.
\end{remark}

\section{Independent, finite-window diagnostic}
We extended the earlier fixed-$X$ orbit sampling to higher starting
altitudes.  For $X=8192$, the actual accelerated map, cutoff $a=786$,
window lower limit $2538$, clipping level $L=3$, and allowances
$K_0=4,K_1=5$, we used $20{,}000$ sampled odd starts for each map and
altitude.  Power $1$ means $[4X,16X)$; powers $2,3,4$ mean
$[X^p,4X^p)$.  All $320{,}000$ sampled trajectories reached the
cutoff.  The first dyadic survival ratio $B_1/B_0$ was:
\begin{center}
\begin{tabular}{@{}lrrrr@{}}
\toprule
$c$ & $[4X,16X)$ & $[X^2,4X^2)$ & $[X^3,4X^3)$ & $[X^4,4X^4)$\\
\midrule
$-1$ & 0.2439 & 0.2564 & 0.2417 & 0.2524\\
$1$ & 0.2535 & 0.2445 & 0.2429 & 0.2511\\
$5$ & 0.3342 & 0.3422 & 0.3548 & 0.3334\\
$13$ & 0.2866 & 0.2805 & 0.2864 & 0.2980\\
\bottomrule
\end{tabular}
\end{center}
These are \emph{finite uniform-band} distributions, not the limiting
logarithmic first-excess entrance measure.  Their stability with
altitude is suggestive but does not give an effective rate for
convergence to $\eta_X$ nor a bound uniform as $X\to\infty$.
This diagnostic uses the standalone executable source included with
this note; it is not a Lean formalisation or an asymptotic proof.

\section{The actual next proof obligation}
The promising new problem is not another parity-cylinder lemma.  It
is a \emph{stopped-law renewal statement}: prove quantitative
non-concentration for the distribution of the next shell-entry state
\emph{after survival conditioning}, together with a control of
ultralong excursions outside the window.  Such a statement would
have to use global first-passage transport or the harmonic predecessor
potential, not merely equidistribution of one fixed finite parity word.
Only after this can the weighted dyadic block-leakage
criterion be applied to establish $T_X=o(X)$.

\begin{thebibliography}{9}
\bibitem{Terras} R. Terras,
\emph{A stopping time problem on the positive integers},
Acta Arith. \textbf{30} (1976), 241--252.
\bibitem{Tao} T. Tao,
\emph{Almost all orbits of the Collatz map attain almost bounded values},
Forum Math. Pi \textbf{10} (2022), e12.
\bibitem{GGM} F. Gon\c{c}alves, R. Greenfeld, J. Madrid,
\emph{Generalized Collatz maps with almost bounded orbits},
arXiv:2111.06170.
\end{thebibliography}
\end{document}
